\section{SHAC: short differentiable windows with a value tail}\label{sec:shac}
\textbf{Xu et al., ICLR 2022.} Accelerated Policy Learning with Parallel
Differentiable Simulation~\cite{shac}.
\loc{Section 3.2, Equations 5--7, Algorithm 1, and Appendix A.1}

\subsection{Reported mechanism in common notation}
SHAC optimizes a policy through short parallel simulator rollouts, adding a
critic bootstrap beyond the window. Windows continue across training updates;
gradients are cut between windows. A reparameterized stochastic policy enables
pathwise differentiation. A target value network supports the actor update,
and a lambda-return critic is fitted from sampled simulator trajectories.
The transition derivatives come from the selected differentiable simulator,
not a learned world model.

With $a_t=\mu_\theta(x_t)+\sigma_\theta(x_t)\odot\epsilon_t$ for a fixed noise
realization, the normalized short-window objective is
\begin{align}
x_{t+1}&=\Fsim(x_t,a_t),\\
L_\pi&=-\frac1{BH}\sum_{b=1}^B\left[
\sum_{t=0}^{H-1}\gamma^t r_t^{(b)}+\gamma^HV_{\bar\psi}(x_H^{(b)})\right],\\
L_V&=\E\big[V_\psi(x_t)-\sg(Y_t)\big]^2,\\
\bar\psi&\leftarrow\rho\bar\psi+(1-\rho)\psi.
\end{align}
The last equation expresses Polyak averaging in parameter form. Terminal
trajectories require special reset and bootstrap handling; concatenating a reset
state into a physical derivative chain is not the intended update.

\subsection{Why the critic changes the gradient problem}
For a one-step window with deterministic action, the action sensitivity includes
\begin{equation}
\nabla_a\big[r(x,a)+\gamma V_\psi(\Fsim(x,a))\big]
=r_a+\gamma F_a^T\nabla_xV_\psi(x').
\end{equation}
The learned value supplies a direction associated with longer-term outcomes
without requiring a long physical Jacobian chain. This shortens differentiation
but introduces dependence on the critic's quality. It does not make difficult
contact branches disappear.

Freezing the target critic's weights still allows its state derivative. Cutting
the simulator chain at a window boundary is different from removing all
derivatives inside the window. Increasing $H$ trades a smaller fraction of
bootstrapped return for a longer, potentially unstable derivative product.
Thus a value tail and horizon selection are coupled design decisions.

\subsection{How to use SHAC as a research control}
A simulator-gradient method answers whether learning the transition function
is helpful for a particular optimization problem. It is not automatically the
correctness oracle: its derivative is that of its numerical contact model,
and a local active-set derivative need not describe a finite change in contact mode.

If a contact-structured learned model beats this baseline, isolate why. Is its
forward prediction better for the evaluated physics? Is the gradient smoother?
Is the model enabling a larger batch or more updates per second? Or does the
critic behave differently? Match these factors before attributing a result to
the model's physical structure.

For a new staircase, known simulator geometry may require no learned dynamics
adaptation, while the policy still needs training. That is a valuable comparison
to the total interaction and compute cost of adapting a learned model. A later
failure of a particular SHAC recipe on G1 tracking should not be generalized
to every SHAC task or simulator.

\section{PWM: pretrained latent dynamics as an optimization interface}\label{sec:pwm}
\textbf{Georgiev et al., ICLR 2025.} PWM: Policy Learning with Multi-Task
World Models~\cite{pwm}.
\loc{Section 3, Equations 4 and 6--10, Algorithm 1, and Appendix C}

\subsection{Reported mechanism in common notation}
PWM pretrains a task-conditioned latent world model and extracts policies by
first-order optimization through imagined rollouts. Its model uses latent
consistency and reward prediction, with SimNorm regularization. Actor training
uses a short return plus a value tail; critics use lambda returns. The method
emphasizes optimization properties of a learned model, rather than prediction
accuracy alone. Optional subsequent online fine-tuning is distinct from the
initial offline policy-extraction phase.

Let $w$ be the learned task embedding, which is not a terrain map. The normalized
state inference and training form is
\begin{align}
z_0&=\mathcal E_\phi(x_0,w),\qquad
\hat z_{k+1}=\mathcal F_\phi(\hat z_k,a_k,w),\\
L_{\rm WM}&=\E_{\mathcal D}\sum_k\gamma^k\left[
\norm{\hat z_{k+1}-\sg(\mathcal E_\phi(x_{k+1},w))}^2
+\operatorname{CE}\big(p_\phi^r(\cdot\mid\hat z_k,a_k,w),T(r_k)\big)\right].
\end{align}
$T$ denotes the log-space categorical reward target. Policy extraction optimizes
Equation~\eqref{eq:shortobjective} with $y=z$ and learned rewards. The paper uses
three value estimators; this is not an ensemble of transition models.

For groups $u^{(i)}$ in a latent vector, SimNorm is
\begin{equation}
\operatorname{SimNorm}(u)=\operatorname{concat}_i
\operatorname{softmax}(u^{(i)}).
\end{equation}
The model uses MLPs with LayerNorm and Mish, with SimNorm on encoder and
transition outputs. These architecture choices belong to PWM; they should not
be silently substituted for another paper's network specification.

\subsection{What regularization can and cannot guarantee}
Each normalized group lies on a probability simplex. Its Jacobian is
\begin{equation}
J_{\rm softmax}=\operatorname{diag}(p)-pp^T.
\end{equation}
This controls the local response of that particular operation and bounds its
output representation. The entire model also contains linear maps and nonlinear
composition, so simplex-valued features alone are not a guarantee of accurate,
small, or useful physical transition sensitivities.

There are at least three separate properties to test: state prediction error,
gradient variability, and the actual return achieved after applying an update.
A low-variance gradient that consistently points toward a wrong optimum is not
a good controller-training signal. Likewise, a larger modeling error can coexist
with a useful finite-scale optimization direction. The synthetic construction
in Section~\ref{sec:common} demonstrates why these quantities are logically separate.

\subsection{Relevance to motion-pretrained transfer}
PWM is the closest conceptual precedent for treating an existing world model as
a resource from which to learn a new controller. But a task embedding can index
training tasks without encoding the spatial geometry of a new staircase. Reusing
a latent model requires testing whether the new interactions remain representable,
not merely attaching a new task ID.

For our question, compare a plain latent/state model and a regularized one with
matched target-data budgets. Hold forward error approximately comparable where
possible, then measure finite-update utility and stability near new contacts.
An improvement would help distinguish beneficial optimization structure from
simply fitting more data with a larger network.

The comparison should include online adaptation. Offline source data may contain
many motions while containing no riser impacts. An actor optimized entirely in
that model can select unsupported actions. A successful extraction result on a
covered task does not quantify the additional contact experience required to
make the new task covered.

\section{DMO: real forward trajectories with learned backward sensitivities}\label{sec:dmo}
\textbf{Amigo et al., CoRL 2025.} First Order Model-Based RL through Decoupled
Backpropagation~\cite{dmo}.
\loc{Section 3.2; Appendix A.1.3, PyTorch decoupling implementation}

\subsection{Reported mechanism in common notation}
DMO generates forward trajectories with a simulator and computes backward
sensitivities through a learned dynamics model evaluated along those trajectories.
Its variants combine this construction with finite-window BPTT, SHAC, or SAPO.
Model fitting uses replay transitions; value targets use simulator trajectories.
The forward/backward separation removes compounding model prediction errors
from the rollout values, while retaining learned transition derivatives.
The paper includes physical Go2 locomotion demonstrations.

Let $x_{t+1}^{\rm sim}$ be a non-differentiable simulator result. A normalized
stop-gradient construction is
\begin{equation}\label{eq:dmo}
\tilde x_{t+1}=\sg(x_{t+1}^{\rm sim})+
\Fhat(\tilde x_t,a_t)-\sg\big(\Fhat(\tilde x_t,a_t)\big).
\end{equation}
Algebraically the forward value is $x_{t+1}^{\rm sim}$; automatic differentiation
assigns the derivatives
\begin{equation}
\frac{\partial\tilde x_{t+1}}{\partial\tilde x_t}=\widehat F_x,
\qquad\frac{\partial\tilde x_{t+1}}{\partial a_t}=\widehat F_a.
\end{equation}
For a Gaussian transition mean/covariance, the normalized maximum-likelihood
fit is the negative log likelihood in Equation~\eqref{eq:nll}. DMO-SHAC then
uses the simulator-valued objective
\begin{equation}
J_H^{\rm DMO}=\E\left[\sum_{t<H}\gamma^t r(x_t^{\rm sim},a_t)
+\gamma^HV_\psi(x_H^{\rm sim})\right]
\end{equation}
with the backward rule in Equation~\eqref{eq:dmo}. DMO-BPTT omits the value tail;
the entropy-regularized variant adds its policy-entropy and critic machinery.

\subsection{Why this is a particularly important control experiment}
Suppose a fully imagined stair rollout diverges after an incorrectly predicted
toe collision. All later states, rewards, and actor observations can then be
evaluated in the wrong contact regime. DMO instead evaluates the next local
model sensitivity at a simulator state. This separates forward distribution
drift from backward model error.

It does not make the learned Jacobians exact. A model can have a poor derivative
at an accurate simulator state, and products of those poor derivatives can
remain misleading. The update is a surrogate-gradient estimator: stop-gradient
algebra does not turn it into the true mathematical derivative of the simulator
return as a function of the actions.

The construction also changes the cost model. A learned rollout can generate
additional hypothetical states without additional simulator interaction. DMO
uses simulator transitions for its forward windows. An adaptation comparison
must count those transitions and their compute time rather than treating the
backward model as a source of free forward data.

\subsection{A diagnostic experiment suggested by this separation}
Use the same learned model, starting states, policy, reward, and horizon in
two updates: one fully imagined and one simulator-forward/model-backward.
Evaluate the resulting policy changes in the simulator. If the decoupled update
works while the imagined one fails, forward rollout drift is implicated. If both
fail despite good one-step predictions, investigate local sensitivities, value
targets, or the task objective.

Neither comparison alone proves that a structured contact model is necessary.
Repeat it with both a plain conditioned model and a structured model. This
crossed design identifies whether structure improves forward accuracy, backward
utility, or both. It directly informs the proposed contact-transfer question
without requiring a differentiable simulator.
